% =============================================================================
%  ANNEXES
% =============================================================================

\chapter{Aide-mémoire du DSL}
\label{annexe:dsl}

\begin{longtable}{@{}p{4.7cm}p{4.4cm}p{4.9cm}@{}}
  \toprule
  \textbf{Jeton} & \textbf{Exemple} & \textbf{Effet} \\
  \midrule
  \endfirsthead
  \toprule
  \textbf{Jeton} & \textbf{Exemple} & \textbf{Effet} \\
  \midrule
  \endhead
  \token{""} & \token{""} & Ne change rien ce temps-là. \\
  \token{|} & \token{|} & Continue l'effet multi-temps précédent. \\
  \token{uni\_n\_R\_G\_B} & \token{uni\_4\_255\_0\_0} &
    Rouge fixe pendant 4 temps. \\
  \token{fond\_n\_c1\_c2} & \token{fond\_8\_255\_0\_0\_0\_0\_255} &
    Fondu rouge $\rightarrow$ bleu sur 8 temps. \\
  \token{str\_n\_R\_G\_B\_f} & \token{str\_2\_255\_255\_0\_120} &
    Strobe jaune pendant 2 temps à 120 Hz. \\
  \token{str\_n\_R\_G\_B} & \token{str\_3\_0\_0\_255} &
    Strobe bleu, fréquence par défaut (60 Hz). \\
  \bottomrule
  \caption{Tous les jetons acceptés.}
\end{longtable}

\paragraph{Règles.} Une composante de couleur doit être comprise entre 0 et
255, sans quoi le jeton est rejeté. La longueur \code{n} est au minimum de 1.
Un préfixe inconnu, un mauvais nombre d'arguments ou un texte non numérique
produisent une \code{DSLError}, journalisée et ignorée pendant la lecture Live.

\paragraph{Fondu sur plusieurs lignes.} On écrit le jeton complet sur la
première ligne, puis un \token{|} sur chacune des lignes suivantes : le nombre
de \token{|} doit valoir \code{n-1} pour couvrir exactement la durée.

\begin{lstlisting}[caption={Fondu de 4 temps (1 ligne d'en-tête + 3 continuations).},numbers=none]
Ligne 1 :  fond_4_255_0_0_0_0_255
Ligne 2 :  |
Ligne 3 :  |
Ligne 4 :  |
\end{lstlisting}

\chapter{Dépannage}
\label{annexe:depannage}

\begin{description}[leftmargin=1.6em,style=nextline]
  \item[Le témoin DMX reste rouge / «{}DMX HS{}.]
    Le serveur ne répond pas : vérifiez qu'\Icare{} tourne encore et que
    l'adresse est bien \code{127.0.0.1:5000}.
  \item[La connexion échoue avec «{}No USB-DMX adapter found{}.]
    Aucun port COM ne contient «{}USB{}» dans sa description. Saisissez le port
    manuellement (par ex. \code{COM3}) dans le champ \emph{Port}. Vérifiez aussi
    que le pilote de l'adaptateur est installé.
  \item[Les projecteurs ne réagissent pas.]
    Contrôlez deux points : le \emph{canal de base} doit correspondre à celui
    programmé sur le projecteur, et la ligne DMX doit être correctement
    chaînée (adaptateur, projecteurs, bouchon de fin si nécessaire).
  \item[L'audio-réactif ne fonctionne pas.]
    Le navigateur doit avoir l'autorisation microphone. L'API
    \code{getUserMedia} n'existe que sur les pages servies en
    \code{localhost} ou en HTTPS : \Icare{} étant local, c'est le cas.
  \item[Le tempo mesuré est bizarre.]
    Les appuis trop espacés (plus de 3 secondes d'écart) sont ignorés.
    Utilisez \textbf{Reset} puis tapez de nouveau, plus régulièrement.
  \item[J'ai un ancien fichier de séquence.]
    Rien à faire : le format historique est migré automatiquement à l'ouverture.
\end{description}

\chapter{Glossaire}
\label{annexe:glossaire}

\begin{description}[leftmargin=1.6em,style=nextline]
  \item[Blueprints] Mécanisme Flask qui regroupe des routes par thème
        (\code{api\_live}, \code{api\_sequence}, \dots).
  \item[Canal de base] Premier canal occupé par un projecteur.
  \item[DMX512] Standard de pilotage de lumières : 512 canaux de 0 à 255 sur un
        bus RS-485 à 250 kbauds.
  \item[DSL] \emph{Domain Specific Language} : ici, le petit langage de chaînes
        qui décrit l'action d'un module sur un temps.
  \item[Enttec USB Pro] Adaptateur USB-DMX très répandu ; son protocole est
        celui implémenté par \file{pydmxIcare.py}.
  \item[Fixture] Un projecteur DMX.
  \item[Frame / image] L'ensemble des 512 valeurs envoyées sur la ligne.
  \item[HSV] Représentation d'une couleur par teinte, saturation et valeur ;
        pratique pour générer des arc-en-ciels.
  \item[Overlay] Effet appliqué par-dessus le mode courant (ici, le strobe).
  \item[RMS] \emph{Root Mean Square} : mesure d'énergie d'un signal sonore.
  \item[Séquence] Fichier JSON décrivant un show ligne par ligne.
  \item[Thread] Fil d'exécution. \Icare{} en utilise deux : la sortie DMX et la
        lecture Live.
  \item[Univers] Les 512 canaux d'une ligne DMX.
\end{description}
